
\conclusion
В ходе выполнения работы был разработан и исследован программный 
комплекс HeatSim3D для численного моделирования нестационарной 
теплопроводности в трёхмерных сплошных средах. Комплекс реализован 
на основе явной конечно-разностной схемы FTCS и включает расчётный 
модуль на языке C++ и графический интерфейс пользователя на языке 
Python.

В ходе работы была реализована постановка задачи с независимыми 
граничными условиями Дирихле на каждой из шести граней расчётной 
области, что позволило моделировать как равномерный нагрев тела, 
так и направленный тепловой поток при различных температурах на 
противоположных гранях. Автоматический выбор шага по времени из 
условия устойчивости исключает возможность получения некорректного 
численного решения при любых допустимых параметрах расчёта.

Проведённая серия вычислительных экспериментов показала, что 
скорость прогрева тела определяется коэффициентом 
температуропроводности материала. Медь и алюминий достигают 
теплового равновесия значительно быстрее конструкционной стали, 
тогда как бетон характеризуется наиболее медленным 
перераспределением тепла. При задании несимметричных граничных 
условий тепловой поток приобретает выраженное направление, что 
наглядно подтверждается визуализацией векторного поля в 
программном комплексе.

Верификация модели путём сравнения численного решения с 
аналитическим показала качественное совпадение характера 
распределения температуры и динамики нестационарного процесса 
теплопроводности, что подтверждает корректность реализованного 
численного алгоритма.

Разработанный программный комплекс может быть использован для 
решения инженерных задач теплопроводности в сплошных средах, 
а также в качестве учебного инструмента для исследования 
нестационарных тепловых процессов. Направлениями дальнейшего 
развития являются реализация неявных численных схем для 
повышения точности расчёта, добавление внутренних источников 
тепла и расширение библиотеки материалов.


Перечень компетенций для направления 09.03.04 «Программная инженерия», сформированных в ходе производственной практики (научно-исследовательской работы):

УК-1. Способен осуществлять поиск, критический анализ и синтез 
информации, применять системный подход для решения задач. 
Компетенция сформирована при изучении литературы по теплопроводности 
и численным методам, решения дифференциальных уравнений с использованием 
электронных библиотек, а также при критическом анализе существующих 
подходов к реализации конечно-разностных схем.

УК-4. Способен осуществлять деловую коммуникацию в устной и 
письменной формах на государственном языке Российской Федерации 
и иностранном(ых) языке(ах). Компетенция проявилась при подготовке 
отчётной документации, оформлении текста научного исследования 
и взаимодействии с руководителем практики. Были использованы нормы 
русского языка и навыки создания научно-технических текстов, в том 
числе при описании математической модели и результатов численного 
моделирования.

ПК-1. Способен проводить научно-исследовательские и 
опытно-конструкторские разработки. Компетенция сформирована при 
разработке численного алгоритма на основе схемы FTCS, реализации 
расчётного модуля на языке C++, проведении вычислительных 
экспериментов и анализе результатов моделирования нестационарной 
теплопроводности.

ПК-4. Способен создавать и анализировать требования на разработку 
программно-информационных систем. Компетенция проявилась при 
формировании требований к программному комплексу HeatSim3D, 
определении входных параметров расчёта, структуры выходных данных, 
требований к точности численного решения и условиям устойчивости 
схемы.

ПК-5. Способен осуществлять концептуальное, функциональное и 
логическое проектирование программно-информационных систем. 
Компетенция проявилась при проектировании архитектуры программного 
комплекса, разработке взаимодействия между расчётным модулем и 
графическим интерфейсом, а также при структурировании логики 
вычислительного процесса и организации средств визуализации 
результатов.